\begin{table}[htbp]
\centering
\caption{SFP sub-period performance (2019--2023). Mirrors Table~7 (DP-PPO sub-periods). Excess CR = SFP minus equal-weight B\&H (same computation as Table~2). SFP outperforms in 3 of 5 calendar years (2020, 2021, 2023). The single largest annual outperformance (+26pp) occurs in 2023, coinciding with the AI investment cycle rally that disproportionately benefited the high-coverage NASDAQ names (NVDA, GOOGL, AVGO) that dominate the SFP basket; it is not possible to rule out that SFP's 2023 excess is a disguised factor exposure rather than a semantic signal effect. SFP underperforms in 2019 and 2022 (rate-hike, drawdown year). Cost = 0.1\% per trade.}
\label{tab:sfp_subperiod}
\begin{tabular}{lrrrr}
\toprule
\textbf{Period} & \textbf{SFP CR (\%)} & \textbf{SFP Sharpe} & \textbf{B\&H CR (\%)} & \textbf{Excess CR (pp)} \\
\midrule
2019 & 35.6 & 1.460 & 43.3 & -7.7 \\
2020 & 61.9 & 1.331 & 51.0 & +11.0 \\
2021 & 46.4 & 1.827 & 42.3 & +4.1 \\
2022 & -30.8 & -0.821 & -29.0 & -1.8 \\
2023 & 81.5 & 2.857 & 55.4 & +26.1 \\
2019--20 pre/bear & 20.4 & 0.595 & 33.8 & -13.5 \\
2020--21 recovery & 179.0 & 2.334 & 139.8 & +39.2 \\
2022--23 rate hike & 23.3 & 0.499 & 9.9 & +13.4 \\
\bottomrule
\end{tabular}
\end{table}
